\section[第14章\q 享受红利的日子]{第14章\\ 享受红利的日子}
\label{sec:ch14}

3年后……

太阳从远处的小山上升起来，光芒闪烁，好像琥珀的颜色。柔伊连续按下3次镜头，拍出完美照片之后，她放下镜头，忘情着看着四周。砖头路上，有一排白色的小屋，蓝色的门窗好像小小的蓝莓。一行红喙海鸥飞入她的视线。她听得到船坞处传来渔船上的绳索轻击的声音。

她已经30岁了吗？想起来，还真不敢相信。这3年过得真快，好像搭了时光机。但另外一面，这3年也发生了许多事。她成了海伦娜的常客，她和乔琪雅成了好朋友。她的人生有了大大的转变。

老妈比医生预料的多活了6个月。柔伊常常过去看望老妈，一些重要的医疗会诊她都参与了。这6个月里母女相聚的时间比这10年以来都更多。亨利称之为“红利”，柔伊觉得非常对。

但是老妈还是过世了，母女甜蜜相聚的时间也就停止了。老爸卖掉了北纽约州的房子，加上保险金，她和老爸买下了布鲁克林区一处温暖公寓。就像亨利说的，拥有了自己的住处。虽然很小，但是父女俩有自己的房间，而柔伊更把自己的小房间变成了工作室，可以改稿，也可以做瑜伽，晚上可以修摄影课。也正如亨利预料的，她用了6个月就存够了摄影课学费。

要付完信用卡债时间就要长一点了。在她把需要支付的款项设置成自动转帐后，她也把信用卡的最低付款金额设置成了自动转帐。不用多付卡债的循环利息，更让她觉得一派轻松。

拿铁因子还真起了作用。

在最初的设置之后，她又照亨利说的，每个月在每张卡上增加了第二笔自动付款金额，定时于付完最低金额后的2周付款。这两张卡就像是能砍断树的锋利斧头，让柔伊终于在22个月后无债一身轻。不用再付利息，拿铁因子又起了作用。

至于学生贷款，这个时间就又更长了。大概要几年时间吧！可是这也没关系的，她会完成的。

光影改变了，琥珀色开始变成了白色。渔船上热闹起来。柔伊举起相机，拍了更多照片。

拿铁因子起作用的地方还不只在信用卡和贷款上，这个写在餐巾纸上的数学公式还非常有效。神奇中的神奇是，在乔琪雅的帮助下，柔伊开始学习烹饪。“就和摄影一样的，”乔琪亚说，“只是在你完成摆盘跟拍照后，你还要吃了它。”柔伊大笑起来，她差点把把拿铁从嘴里喷出来。不过，因为自己准备午餐，可以省下的钱，就可不是玩笑的。就像拜伦戒了烟，柔伊一边想着，一边微笑了。

柔伊还停掉了从来都不看的电视节目，停掉了她从来也不去运动的健身房会员。处理了她从来也没空穿的衣服，甚至还丢弃了时装目录，免得受诱惑。同时，她的退休帐户和旅游帐户里，都开始有钱潮涌入。

太阳稍微升高，小村因此更加热闹起来，她听着渔民在准备出海时低声交谈。一天最黄金的时间就要来了。

她拍了一张又一张的照片，然后，开始研究手上的相机。真是个漂亮的小物。那是，拜伦和乔琪雅送的生日礼物，就在她要开始旅行之前，他们提前送的。

而今天正是她的生日。

这是她希腊岛6周游的最后一天。

这段时间，她每天都写点东西，昨天刚刚把全部内容寄出去。芭芭拉回信说，杂志社内的同事想把她的文章和照片做成一个专栏。她回去之后就晋升了：柔伊现在不只是位编辑，还是一位专栏作家。芭芭拉在邮件尾处署名：

{\kaishu 生日快乐，小伊。}

{\kaishu 你家老大}

{\kaishu 顺便说一下，明信片很棒。}

柔伊每年都出游一趟，已经连续第三年了，她高调的休假研习。去年她去密西西比河，用了5周的时间，在美国西南的山脉，从、亚利桑那州的圣多娜到新墨西哥州的拉斯克鲁赛斯。她的有些照片，还挂到了海伦娜咖啡厅的墙上。

第二年的旅游非常棒，但第一年的圣多娜红石风景却是无法超越的。那年秋天，老妈走了。老爸和她花了4周的时间去缅因州。他们摘野莓吃，拍秃鹫，搭龙虾船出海。父女俩聊着老妈的过去，交换了他们心目中的珍贵画面。这趟旅行费用不贵，因为柔伊还要存去希腊的旅游款，但是却非常丰富。

在要来希腊之前，她问老爸，如果让他选，他最想去哪里？“阿拉斯加。”老爸毫不犹豫的回答。

“那就说定了。”柔伊说，“明年。老爸。可以开始准备了。”

柔伊喝了一口希腊咖啡，看着小渔村在希腊的阳光底下热闹起来。她举起相机，靠近了眼睛，通过了相机眼窗，看着世界，并按下快门。

“太棒了！”她低声说。

柔伊‧丹尼尔斯今年30岁。她是世上最富足的女人，这点，她比谁都清楚。
